\section{习题课 \#11（11月22日）}

\subsection{正规矩阵与谱定理}

在复数域上，用 $A^*$ 表示共轭转置。

\begin{definition}
矩阵 $A\in\C^{n\times n}$ 称为正规矩阵，如果
\begin{equation}
 A^*A=AA^*.
\end{equation}
Hermite 矩阵、酉矩阵和复数域上的对角矩阵都是正规矩阵。
\end{definition}

\begin{theorem}[正规矩阵谱定理]
矩阵 $A\in\C^{n\times n}$ 正规，当且仅当存在酉矩阵 $U$ 与对角矩阵 $D$ 使
\begin{equation}
 A=UDU^*.
\end{equation}
\end{theorem}

\begin{proof}
若 $A=UDU^*$，则
\begin{equation}
 A^*A=U D^*D U^*=U DD^*U^*=AA^*.
\end{equation}

反之，由 Schur 分解，存在酉矩阵 $U$ 使
\begin{equation}
 T=U^*AU
\end{equation}
为上三角矩阵。正规性在酉相似下保持，所以 $T$ 正规。比较 $T^*T$ 与 $TT^*$ 的第一个对角元可知第一行非对角元素全为零，再对右下角子块归纳，得到 $T$ 必为对角矩阵。
\end{proof}

\begin{corollary}[Hermite 与实对称谱定理]
\begin{enumerate}
  \item 若 $A=A^*$，则 $A$ 的所有特征值均为实数，且可酉对角化；
  \item 若 $A\in\R^{n\times n}$ 且 $A=A^T$，则存在实正交矩阵 $Q$ 使
  \begin{equation}
   A=QDQ^T,
  \end{equation}
  其中 $D$ 为实对角矩阵。
\end{enumerate}
\end{corollary}

\begin{theorem}[不同特征值的特征向量正交]
若 $A$ 正规，且
\begin{equation}
 Ax=\lambda x,
 \qquad
 Ay=\mu y,
 \qquad
 \lambda\neq\mu,
\end{equation}
则
\begin{equation}
 x^*y=0.
\end{equation}
\end{theorem}

\begin{proof}
谱定理给出一组标准正交特征向量基，不同特征值对应的特征空间由互不相交的基向量子集张成，故彼此正交。也可利用正规矩阵满足
\begin{equation}
 \Ker(A-\lambda I)=\Ker(A^*-\overline\lambda I)
\end{equation}
直接计算。
\end{proof}

\subsection{实矩阵在实数域与复数域上的相似性}

\begin{theorem}
设 $A,B\in\R^{n\times n}$。则 $A$ 与 $B$ 在 $\R$ 上相似，当且仅当它们在 $\C$ 上相似。
\end{theorem}

\begin{proof}
实相似显然蕴含复相似。反之，假设存在复可逆矩阵
\begin{equation}
 S=P+\ii Q,
 \qquad P,Q\in\R^{n\times n},
\end{equation}
使
\begin{equation}
 AS=SB.
\end{equation}
比较实部与虚部得到
\begin{equation}
 AP=PB,
 \qquad
 AQ=QB.
\end{equation}
因此对每个 $t\in\R$，
\begin{equation}
 A(P+tQ)=(P+tQ)B.
\end{equation}
多项式 $p(t)=\det(P+tQ)$ 不恒为零，因为 $p(\ii)=\det(P+\ii Q)\neq0$。它在实数轴上只有有限多个零点，所以可选 $t\in\R$ 使 $P+tQ$ 可逆，从而得到实相似变换。
\end{proof}

\subsection{由一个矩阵生成的多项式代数}

\begin{definition}
对 $A\in\F^{n\times n}$，定义
\begin{equation}
 \F[A]=\set{f(A):f\in\F[t]},
\end{equation}
并定义 $A$ 的中心化子
\begin{equation}
 \mathcal C(A)=\set{B\in\F^{n\times n}:AB=BA}.
\end{equation}
\end{definition}

\begin{proposition}
$\F[A]$ 与 $\mathcal C(A)$ 都是 $\F^{n\times n}$ 的线性子空间，并且
\begin{equation}
 \F[A]\subseteq\mathcal C(A).
\end{equation}
若 $B\in\mathcal C(A)$，则
\begin{equation}
 \F[B]\subseteq\mathcal C(A).
\end{equation}
\end{proposition}

\begin{proof}
若 $AB=BA$，则归纳可得 $AB^k=B^kA$，继而 $A$ 与任意 $B$ 的多项式可交换。取 $B=A$ 即得第一项包含关系。
\end{proof}

\begin{proposition}[正定平方根是多项式]
设 $B\in\R^{n\times n}$ 为实对称正定矩阵。则存在多项式 $p$，使
\begin{equation}
 B^{1/2}=p(B),
\end{equation}
其中 $B^{1/2}$ 是唯一的实对称正定平方根。因此，若 $AB=BA$，则
\begin{equation}
 AB^{1/2}=B^{1/2}A.
\end{equation}
同理，$B^{-1/2}$ 也可写成 $B$ 的多项式。
\end{proposition}

\begin{proof}
设 $B$ 的互异特征值为 $\lambda_1,\ldots,\lambda_s>0$。由 Lagrange 插值，可取多项式 $p$ 满足
\begin{equation}
 p(\lambda_j)=\sqrt{\lambda_j},
 \qquad j=1,\ldots,s.
\end{equation}
若 $B=Q\diag(\lambda_1',\ldots,\lambda_n')Q^T$，则
\begin{equation}
 p(B)=Q\diag(\sqrt{\lambda_1'},\ldots,\sqrt{\lambda_n'})Q^T=B^{1/2}.
\end{equation}
\end{proof}

\begin{remark}
矩阵的三个常见分类关系分别为
\begin{equation}
 B=PAQ\quad\text{（等价）},
\end{equation}
\begin{equation}
 B=P^{-1}AP\quad\text{（相似）},
\end{equation}
以及
\begin{equation}
 B=P^TAP\quad\text{（合同）}.
\end{equation}
它们分别对应一般线性映射、线性变换和双线性型或二次型的换基。典型标准形分别由秩、Jordan 或有理标准形、以及惯性指数描述。
\end{remark}

\subsection{Fibonacci 数列的矩阵解法}

设
\begin{equation}
 a_0=0,
 \qquad a_1=1,
 \qquad a_{n+2}=a_{n+1}+a_n.
\end{equation}
令
\begin{equation}
 \alpha_n=\begin{pmatrix}a_{n+1}\\a_n\end{pmatrix},
 \qquad
 M=\begin{pmatrix}1&1\\1&0\end{pmatrix}.
\end{equation}
则
\begin{equation}
 \alpha_{n+1}=M\alpha_n,
 \qquad
 \alpha_n=M^n\alpha_0.
\end{equation}

矩阵 $M$ 的特征值为
\begin{equation}
 \varphi=\frac{1+\sqrt5}{2},
 \qquad
 \psi=\frac{1-\sqrt5}{2}=-\varphi^{-1}.
\end{equation}
对角化后得到 Binet 公式
\begin{equation}
 a_n=\frac{\varphi^n-\psi^n}{\sqrt5}.
\end{equation}
该方法同时说明递推数列、矩阵幂和特征值之间的直接联系。

\subsection{诱导矩阵二范数与谱半径}

\begin{definition}
对 $A\in\C^{m\times n}$，诱导二范数定义为
\begin{equation}
 \norm{A}_2
 =\sup_{x\neq0}\frac{\norm{Ax}_2}{\norm x_2}.
\end{equation}
对方阵 $A$，谱半径定义为
\begin{equation}
 \rho(A)=\max_{\lambda\in\sigma(A)}|\lambda|.
\end{equation}
\end{definition}

\begin{proposition}
诱导二范数满足：
\begin{enumerate}
  \item $0\leq\norm A_2<\infty$，且 $\norm A_2=0$ 当且仅当 $A=0$；
  \item $\norm{aA}_2=|a|\norm A_2$；
  \item $\norm{A+B}_2\leq\norm A_2+\norm B_2$；
  \item $\norm{AB}_2\leq\norm A_2\norm B_2$；
  \item $\norm A_2^2=\rho(A^*A)$。
\end{enumerate}
\end{proposition}

\begin{theorem}
对任意方阵，
\begin{equation}
 \rho(A)\leq\norm A_2.
\end{equation}
若 $A$ 正规，则
\begin{equation}
 \rho(A)=\norm A_2.
\end{equation}
\end{theorem}

\begin{proof}
若 $Av=\lambda v$，则
\begin{equation}
 |\lambda|\norm v_2=\norm{Av}_2\leq\norm A_2\norm v_2.
\end{equation}
对正规矩阵写成 $A=UDU^*$，酉不变性给出
\begin{equation}
 \norm A_2=\norm D_2=\max_i|d_i|=\rho(A).
\end{equation}
\end{proof}

\begin{remark}
一般矩阵中不一定取等号。例如
\begin{equation}
 N=\begin{pmatrix}0&1\\0&0\end{pmatrix}
\end{equation}
满足 $\rho(N)=0$，但 $\norm N_2=1$。
\end{remark}

\subsection{矩阵指数与谱半径为零的判别}

由次乘法性，级数
\begin{equation}
 \e^A=\sum_{k=0}^{\infty}\frac{A^k}{k!}
\end{equation}
绝对收敛，因为
\begin{equation}
 \sum_{k=0}^{\infty}\frac{\norm A_2^k}{k!}
 \leq\e^{\norm A_2}.
\end{equation}
因此矩阵指数定义良好，并满足 $AB=BA$ 时
\begin{equation}
 \e^{A+B}=\e^A\e^B.
\end{equation}

\begin{theorem}
对 $A\in\C^{n\times n}$，下列条件等价：
\begin{enumerate}
  \item $\rho(A)=0$；
  \item $A$ 的全部特征值为 $0$；
  \item $A$ 幂零；
  \item 对每个 $k\geq1$，$\tr(A^k)=0$。
\end{enumerate}
\end{theorem}

\begin{proof}
前三项由 Cayley--Hamilton 定理直接等价。若特征值全为零，则 $\tr(A^k)$ 是特征值 $k$ 次幂之和，故为零。

反之，设互异非零特征值为 $\lambda_1,\ldots,\lambda_s$，其代数重数为 $m_1,\ldots,m_s$。由
\begin{equation}
 \sum_{j=1}^s m_j\lambda_j^k=0,
 \qquad k=1,\ldots,s,
\end{equation}
除去一个 $\lambda_j$ 后得到 Vandermonde 型线性方程组。其系数矩阵可逆，迫使 $m_j\lambda_j=0$，矛盾。因此不存在非零特征值。
\end{proof}

\subsection{半正定矩阵的等价刻画}

\begin{theorem}
对实方阵 $A$，下列条件等价：
\begin{enumerate}
  \item 存在实矩阵 $B$ 使
  \begin{equation}
   A=B^TB;
  \end{equation}
  \item $A=A^T$ 且
  \begin{equation}
   x^TAx\geq0
   \qquad\text{对所有 }x\in\R^n;
  \end{equation}
  \item 存在唯一的实对称半正定矩阵 $C$ 使
  \begin{equation}
   A=C^2.
  \end{equation}
\end{enumerate}
\end{theorem}

\begin{proof}
若 $A=B^TB$，则 $A$ 对称且
\begin{equation}
 x^TAx=\norm{Bx}_2^2\geq0.
\end{equation}
若 $A$ 对称半正定，谱定理给出
\begin{equation}
 A=Q\diag(\lambda_1,\ldots,\lambda_n)Q^T,
 \qquad \lambda_i\geq0.
\end{equation}
取
\begin{equation}
 C=Q\diag(\sqrt{\lambda_1},\ldots,\sqrt{\lambda_n})Q^T
\end{equation}
即可。半正定平方根的唯一性可在每个特征空间上验证。最后令 $B=C$ 即回到第一项。
\end{proof}
